\section{Limitations and threats to validity}\label{sec:limits}
\begin{itemize}
  \item All evidence is synthetic or trace-driven. There is no production workload, and the production-flow scenario was
        written by the author.
  \item One cloud provider and region, two instance types, one placement group (a low-latency network).
  \item The v2 budgets were set after seeing earlier offline data. Trace simulations reconstruct arrivals where the
        trace lacks them, and Huawei timestamps are completion times.
  \item Durations are learned from successful handlers only (success-conditioned labels).
  \item The broker is a prototype: in-memory, no durability, authentication or backpressure, and the lowest throughput
        of the three brokers tested. E2 has two repetitions, so its checks have no intervals.
  \item Shortest-expected-first with a non-preemptive scheduler cannot shorten a long job's own service time; we make no
        tail-latency claim, and long jobs can lose (Section~\ref{sec:results}).
  \item Alibaba GPU traces and Kafka, Pulsar and managed queues are out of scope.
\end{itemize}
\paragraph{Reproducibility.} Code, configurations, pre-registrations and result files are in the repository
(\url{https://github.com/divyanshusingh2903/predictive-multi-level-message-queue}); raw per-message data is archived
outside git with SHA-256 sums. \todo{record the release tag and archive identifier.}

\paragraph{AI assistance.} \todo{Disclose AI-assisted drafting per the venue's current policy (issue \#27).}
